\section{Theory}\label{sec:theory}

Throughout, statistics are bounded, $Z_k\in[z_{\mathrm{lo}},z_{\mathrm{hi}}]$ with width $w\defeq z_{\mathrm{hi}}-z_{\mathrm{lo}}$, and bets and references are predictable. No independence or exchangeability is assumed.

\subsection{No false revocation}

\begin{theorem}[Uniform validity]\label{thm:validity}
Fix $\alpha\in(0,1)$. If horizon $h$ satisfies \eqref{eq:null-mean}, then $\PR(\exists K:\ W^{(h)}_K\ge1/\alpha)\le\alpha$,
uniformly over all predictable planners and betting rules. Running the $H$ horizons at level $\alpha/H$ each gives
$\PR(\exists t,\ \exists h\in\Hcal_0:\ h\text{ revoked by round }t)\le\alpha$ for the set $\Hcal_0$ of faithful horizons.
\end{theorem}
\begin{proof}
By Proposition~\ref{prop:evar} and linearity, $W^{(h)}$ is a nonnegative supermartingale with $W^{(h)}_0=1$ under the
null, using only \eqref{eq:null-mean} and predictability. Ville's inequality bounds $\PR(\sup_KW^{(h)}_K\ge H/\alpha)\le\alpha/H$ for each horizon, and a union
bound over $|\Hcal_0|\le H$ horizons gives the simultaneous statement.
\end{proof}

The guarantee is conservative: under the null the wealth is a supermartingale that tends to drift down, so empirical rates
sit far below $\alpha$ (Section~\ref{sec:exp-synth}). It controls false \emph{revocations}; it says nothing about trusting an
unfaithful horizon, which is the job of the certificate below.

\subsection{A confidence sequence for the average conditional mean}

The simulation-lemma-style bound needs the running average $m_K=\frac1K\sum_{k\le K}\mu_k$ of the \emph{conditional} means, which vary
over the deployment. A betting confidence sequence for a \emph{constant} mean does not certify it: its proof tests
$\E[Z_k\mid\Fcal_{k-1}]\ge u$ for every $k$, which is not implied by an average bound. We use an
empirical-Bernstein supermartingale instead.

\begin{lemma}[Confidence sequence for the average conditional mean]\label{lem:cs}
Let $\widehat Z_k\in[z_{\mathrm{lo}},z_{\mathrm{hi}}]$ be predictable, $\psi(\lambda)=-\log(1-\lambda)-\lambda$,
$V_K=\sum_{k\le K}\big((Z_k-\widehat Z_k)/w\big)^2$, and fix a grid $\lambda_1,\dots,\lambda_n\in(0,1)$ with
prior weights $\nu_j$. With probability at least $1-\alpha$, simultaneously for all $K\ge1$ and all $j$,
\begin{equation}\label{eq:cs}
m_K\ \le\ U_K\defeq\bar Z_K+\frac{w}{\lambda_jK}\Big(\psi(\lambda_j)V_K+\log\frac1{\alpha\nu_j}\Big),
\end{equation}
so the minimum over $j$ is a valid upper confidence sequence. The width is of the empirical-Bernstein order
$\sqrt{\hat\sigma_K^2\log(1/\alpha)/K}+w\log(1/\alpha)/K$.
\end{lemma}
\begin{proof}
Rescale to $[0,1]$. For $y\ge-1$ and $\lambda\in[0,1)$, $\exp(\lambda y-\psi(\lambda)y^2)\le1+\lambda y$
\citep{fan2015exponential}. Put $y=\widehat x_k-x_k\in[-1,1]$. Then
$\E[\exp(\lambda(\mu_k-x_k)-\psi(\lambda)(x_k-\widehat x_k)^2)\mid\Fcal_{k-1}]\le e^{\lambda(\mu_k-\widehat x_k)}(1+\lambda(\widehat x_k-\mu_k))\le1$,
for \emph{any} value of the conditional mean $\mu_k$. So $M_K(\lambda)=\exp(\lambda\sum_k(\mu_k-x_k)-\psi(\lambda)V_K)$ is a nonnegative supermartingale, and so is
$\sum_j\nu_jM_K(\lambda_j)$. Since $\nu_jM_K(\lambda_j)$ is dominated by the mixture, Ville's inequality gives
$\PR(\exists K,j:\ \nu_jM_K(\lambda_j)\ge1/\alpha)\le\alpha$; rearranging gives \eqref{eq:cs}.
\end{proof}

\begin{corollary}[Windows]\label{cor:window}
Restart the sequence at block starts $s_i=ib$ and give start $s_i$ the level $\alpha/((i+1)(i+2))$. Then with probability at least $1-\alpha$, simultaneously
for all starts and all end times, the average conditional mean over $(s_i,t]$ is at most the corresponding bound.
Reading the window with the latest start $s_i\le t-w_0$ gives an anytime-valid bound on the \emph{recent} average over a window of length in $[w_0,w_0+b)$ (for $t<w_0$ it is the whole history).
\end{corollary}
\begin{proof}
Apply Lemma~\ref{lem:cs} at each start with level $\alpha_i=\alpha/((i+1)(i+2))$ and use $\sum_i\alpha_i\le\alpha$.
\end{proof}

The windowed form matters: the full-history average moves slowly after a change, so a certificate meant to follow the
present must be read over a window; the price is the larger $\log(1/\alpha_i)\approx2\log i$ term.

\subsection{The certified horizon and a simulation-lemma consequence}

\begin{theorem}[Anytime-valid certified horizon]\label{thm:cert}
Suppose each horizon's statistic is certifiable with constants $(c,\bar\rho)$ (Definition~\ref{def:cert}) and
$\rho=\TV$ with modulus $\phi$. Apply Lemma~\ref{lem:cs} (or Corollary~\ref{cor:window}) to the chain of each horizon
$h\le H$ at level $\alpha/H$, giving bounds $U^{(h)}_t$ for the average over a chain index set $\mathcal S_{h,t}$ (full history or a window).
Then with probability at least $1-\alpha$, simultaneously for all rounds $t$ and horizons $h\le H$,
\begin{equation}\label{eq:certtv}
\frac1{|\mathcal S_{h,t}|}\sum_{k\in\mathcal S_{h,t}}\TV\big(Q_{s_k,h},P^M_{s_k,h}\big)\ \le\ \min\big\{1,\ \phi\big(cU^{(h)}_t+\bar\rho\big)\big\}.
\end{equation}
Consequently, with probability at least $1-\alpha$, for all $t$ and all $h\le\hcert_\tau(t)$ the left side is at most $\tau$. The statement is about the \emph{window average} over logged issue times: a short spike inside a long window can be diluted below $\tau$, and it concerns the logged action sequences, not a counterfactual imagined policy.
\end{theorem}
\begin{proof}
By certifiability the average divergence over $\mathcal S_{h,t}$ is at most $c\,m+\bar\rho$ with $m$ the average conditional mean, which Lemma~\ref{lem:cs} bounds by
$U^{(h)}_t$ for all $t,h$ with probability $1-\alpha$ after a union bound over $H$ horizons. By the modulus, concavity of $\phi$ (Jensen) and monotonicity,
the average of $\TV\le\phi(\dS)$ is at most $\phi(c\,m+\bar\rho)\le\phi(cU^{(h)}_t+\bar\rho)$; total variation is at most one.
\end{proof}

\begin{corollary}[Value gap]\label{cor:value}
Let rewards $r\in[0,R_{\max}]$ be a function of the outcome and let actions be held fixed (open-loop), as in the action-conditioned predictive. For discount weights $w_i\ge0$, on the event of Theorem~\ref{thm:cert}
the weighted sum of window-averaged per-step reward gaps satisfies
$\sum_{i\le h}w_i\,G_i(t)\le R_{\max}\sum_{i\le h}w_i\min\{1,\phi(cU^{(i)}_t+\bar\rho)\}$, where $G_i(t)$ is the average of $|\E_M r(Y^{(i)})-\E_Q r(Y^{(i)})|$ over horizon $i$'s chain index set.
\end{corollary}
\begin{proof}
For fixed actions $|\E_Pr-\E_Qr|\le R_{\max}\TV(P,Q)$; apply \eqref{eq:certtv} for each $i$.
\end{proof}

\begin{remark}[Relation to the classical simulation lemma]
The classical lemma \citep{kearns2002near} converts a \emph{known} one-step error into an $h$-step value gap, paying a compounding factor such as $\gamma/(1-\gamma)^2$.
Here each horizon's own chain is certified, so there is no propagation step and no concentrability coefficient (the price is that the certified quantity is the discrepancy on the logged action sequences, not on counterfactual rollouts of a new policy); the price is that horizon $h$ is certified from a $1/h$ share
of the rounds. When the gap profile is constant across issue times in the window, $\sum_iw_iG_i$ is the average discounted return gap. We are not aware of a
time-uniform version of this statement in the literature.
\end{remark}

\subsection{Detection delay}

Take the e-detector $R_t$ built from the mixture bet. For the alarm time $T=\inf\{t:R_t\ge1/\alpha\}$ the classical argument gives the false-alarm guarantee, and a martingale
concentration argument gives a finite-sample delay.

\begin{theorem}[Finite-sample detection delay]\label{thm:delay}
Let $X_k=Z_k-\eps$ (so $|X_k|\le w$, the width of the statistic's range). \emph{(a)} Under the null, $\E_\infty[T]\ge1/\alpha$. \emph{(b)} Suppose that after a changepoint $\tau$,
$\E[X_k\mid\Fcal_{k-1}]\ge\Delta>0$ and $\E[X_k^2\mid\Fcal_{k-1}]\le v$ for all $k>\tau$, and that the grid is $\lambda_j=j/(mw)$, $j=1,\dots,m$ (the one the implementation uses) and
$\Delta\ge 2v/(mw)$. Let $\lambda^\star=\min\{\Delta/(2v),\,1/(2w)\}$. For every $\delta\in(0,1)$ and
\begin{equation}\label{eq:N}
n\ \ge\ N\defeq\max\Big\{\,\frac{8\log(m/\alpha)}{\lambda^\star\Delta},\ \frac{(80v+8w\Delta)\log(1/\delta)}{\Delta^2}\Big\},
\end{equation}
we have, for every history, $\PR(T-\tau>n\mid\Fcal_\tau,\,T>\tau)\le\delta$. Note $8/(\lambda^\star\Delta)\le16(v/\Delta^2+w/\Delta)$.
\emph{(c)} For horizon $h$ on a non-overlapping chain, with per-chain margin $\Delta_h$ and second moment $v_h$, the delay in rounds is at most $h(N_h+1)$.
\end{theorem}

The proof is in Appendix~\ref{app:delay}. When the first term of \eqref{eq:N} dominates and the variance term is the larger one ($v\gg w\Delta$), $N\approx16\,v\log(m/\alpha)/\Delta^2$, which is within a constant factor of the Lorden--Lai rate $2\sigma^2\log(1/\alpha)/\Delta^2$ for a Gaussian-like pair with the same mean shift and variance \citep{lorden1971,lai1998information} (the comparison is of a $1-\delta$ quantile with an expected delay, so it is a statement about order, not about constants; the second term of \eqref{eq:N} dominates for $\delta\lesssim0.2$). The bound is loose in practice: on the simulator of Section~\ref{sec:exp-horizon} it gives $N$ in the thousands for horizon one against a measured delay near $500$. When the
bet saturates ($v\lesssim w\Delta$) the delay scales as $w/\Delta$ instead of $1/\Delta^2$. Part (c) is the horizon-resolved law: horizon $h$ pays a factor $h$ for fewer updates and
gains through a larger margin $\Delta_h$ as the shift compounds along the rollout. Section~\ref{sec:exp-horizon} tests the first-order prediction $h(\log(H/\alpha)+\log m)/D^\star_h$ with $D^\star_h=\sup_\lambda\E\log(1+\lambda Z)$.

\begin{remark}[The e-process alone]
For the plain e-process $W$ started at $1$, the same argument holds with $\log(m/\alpha)$ replaced by $\log(m/(\alpha W_\tau))$, where $W_\tau$ is the wealth at the changepoint: wealth lost before the change
delays the alarm. The Shiryaev--Roberts form removes this term because it contains the term started at $\tau{+}1$ at full wealth.
\end{remark}
